\documentclass[10pt,a4paper]{amsart}
\usepackage[margin=19mm]{geometry}
\usepackage{amsmath,amssymb,mathtools}
\usepackage[scr=rsfs,cal=euler]{mathalfa}
\usepackage{xfrac}
\usepackage[unicode,hidelinks,bookmarks=false]{hyperref}
\newcommand{\ie}{i.e.\ }
\newcommand{\cf}{cf.\ }
\newcommand{\resp}{respectively\ }
\newcommand{\Subsection}{\subsection}
\providecommand{\nobreakdash}{\nobreak\hbox{-}}
\setlength{\emergencystretch}{2em}
\multlinegap0pt
\usepackage{bm}
\usepackage{amssymb}
\usepackage{enumerate}
\usepackage[all]{xy}
\newtheorem{proposition}{Proposition}
\newtheorem{remark}{Remark}
\newtheorem{fact}{Fact}
\newtheorem{lemma}{Lemma}
\newtheorem{claim}{Claim}
\title{Complete Characterization on Maximum Pairwise Cross Intersecting Families (I)}
\author{Yang Huang, Yuejian Peng}
\date{Source: arXiv:2601.01929v1 (CC BY 4.0)}
\begin{document}
\maketitle

\noindent\emph{Source and licence.} Complete Characterization on Maximum Pairwise Cross Intersecting Families (I). Version arXiv:2601.01929v1 (CC BY 4.0), distributed by arXiv under the Creative Commons Attribution 4.0 International licence, http://creativecommons.org/licenses/by/4.0/, as recorded in the arXiv licence metadata for this exact version and shown on the abstract page https://arxiv.org/abs/2601.01929v1. Full licence notice: \texttt{licenses/arxiv-2601.01929-CC-BY-4.0.txt}.\par\smallskip

\noindent\textit{Editorial scope.} Counted unit: the article's proposition 7-18-1 (source0.tex lines 868--876). The article's complete original proof of it is delivered with it (source0.tex lines 878--910, one \texttt{proof} environment of 33 lines). No mathematics is written, repaired or re-derived: the statement and the proof are the article's own lines, in the article's own notation, and no retained line is substituted or re-flowed.  Retained uncounted as the source-local context the proof needs: lines 419--421; lines 490--492; lines 494--499; lines 555--557; lines 566--568; lines 589--604; lines 777--780; lines 787--791.  Nothing was omitted from any retained span, so the package carries no omission record.  No retained line contains a citation, so the wrapper carries no bibliography and no bibliography entry.  No retained line contains a figure environment, an image-inclusion command or drawing code, so no image is delivered and the package declares no asset dependency.  Editorial formatting only, in addition: an AMS \texttt{amsart} preamble that restates the article's own declarations and macros used by the retained lines, namely the environments \texttt{proposition}, \texttt{remark}, \texttt{fact}, \texttt{lemma}, \texttt{claim} on separate counters, together with the standard packages those lines need.  The retained environments print in this document as Proposition 1, Remark 1, Fact 1, Lemma 1, Claim 1 in order of appearance, whereas the article prints the same items with the numbering of its own full document.  The complete native source of the article as distributed in the arXiv e-print for this exact version is bundled unchanged as \texttt{sources/192121/source0.tex}.
\begin{proposition}\label{p2.7}
Let $a, b, n$ be positive integers and $a+b\leq n$. For $P\subseteq [n]$ with $|P|\leq a$, let $Q$ be the partner of $P$. Then $\mathcal{L}(Q, b)$ is the maximum  $b$-uniform family that is cross intersecting to $\mathcal{L}(P, a)$.
\end{proposition}

\begin{remark}\label{r2.6}
Let $F\subseteq [n]$ with $|F|=f$ and $k\leq n-f$. Suppose that  $H$ is the partner of $F$, and $K$ is the k-partner of $F$, then we have $\mathcal{L}(H, k)=\mathcal{L}(K, k)$.
\end{remark}

By Proposition \ref{p2.7} and Remark \ref{r2.6}, we have the following fact.
\begin{fact}\label{f2.8}
Let $a, b, n$ be positive integers and $n\geq a+b$. For  $A\subseteq [n]$ with $|A|=a$, let $B$ be the $b$-partner of $A$, then
 $\mathcal{L}(B, b)$ and $\mathcal{L}(A, a)$ are cross intersecting, moreover,
  $\mathcal{L}(B, b)$ is maximal to $\mathcal{L}(A, a)$.
\end{fact}

\begin{fact}\label{f2.10}
Let $a, b, n$ be positive integers and $n\geq a+b$. Suppose that $A$ is an $a$-subset of $[n]$, and $B$ is the $b$-partner of $A$. Let $A'$ be the  $a$-partner of $B$, then $(A', B)$ is maximal. Moreover, $A\prec A'$.
\end{fact}

\begin{fact}\label{fact2.7}
Let $a, b, k, n$ be positive integers and $n\geq \max\{a+k, b+k\}$. Let $A$  and $B$ be two subsets of $[n]$ with sizes $|A|=a$ and $|B|=b$. Suppose that $K_a$ and $K_b$ are the $k$-partners of $A$ and $B$ respectively. If $A\prec B$, then $K_b\prec K_a$. In particular, if $A$ is the $a$-parity of $B$, then $K_a= K_b$.
\end{fact}

\begin{proposition}\label{prop2.9}
Let $f, g, h, n$ be positive integers with $f\geq g$ and $n\geq f+h$. Let
\begin{align*}
&\mathcal{F}=\left\{ F\in{[n]\choose f}: \text{ there exists } H\in {[n]\choose h} \text{ such that } (F, H) \text{ is maximal }  \right\},\\
&\mathcal{G}=\left\{ G\in{[n]\choose g}: \text{ there exists } H\in {[n]\choose h} \text{ such that } (G, H) \text{ is maximal }  \right\}.
\end{align*}
Let $\mathcal{H}_{\mathcal{F}}\subseteq {[n]\choose h}$ and $\mathcal{H}_{\mathcal{G}}\subseteq {[n]\choose h}$ be the families
such that $\mathcal{F}$ and $\mathcal{H}_{\mathcal{F}}$ are maximal pair families and $\mathcal{G}$ and $\mathcal{H}_{\mathcal{G}}$ are  maximal pair families.
Then the following properties hold:

(i) $\mathcal{F}$ is the $f$-parity of $\mathcal{G}$;

(ii) $\mathcal{H}_{\mathcal{G}}\subseteq \mathcal{H}_{\mathcal{F}}$;

(iii) for any $G\in \mathcal{G}$, let $F\in \mathcal{F}$ be the $f$-parity of $G$ and let $H\in \mathcal{H}$ be such that $(F, H)$ is maximal, then $(G, H)$ is maximal.
\end{proposition}

\begin{remark}\label{remark2.20+}
Let $d\geq 2$ be an integer. We consider $d$ L-initial families $\mathcal{L}([n], A_1, a_1)$, $\dots$, $\mathcal{L}([n], A_d, a_d)$ with $|A_1|=a_1, \dots, |A_d|=a_d$. For some fixed $i\in [d]$, let $S\subseteq [d]\setminus \{i\}$ be the set of all $j\in [d]\setminus \{i\}$ satisfying that $n\geq a_i+a_j$. Suppose that $\{1, n-a_i+2, \dots, n\}\preceq  A_i$. If $\mathcal{L}(A_i, a_i)$ and $\mathcal{L}(A_j, a_j)$ are cross intersecting for each $j\in S$, then
$\mathcal{L}(A_j, a_j)$, $j\in S$ are pairwise cross intersecting families since for any $j\in S$, every member of $\mathcal{L}(A_j, a_j)$ contains $1$.
\end{remark}

\begin{lemma}\label{>}
Let $k_1\geq k_2\geq\dots\geq k_t\geq 2$ and $k_1+k_3\leq n <k_1+k_2$.
Suppose that $(\mathcal{F}_1, \dots, \mathcal{F}_t)$ is an L-initial extremal $(n, k_1, \dots, k_t)$-cross intersecting system with IDs $I_1, I_2, \dots, I_t$  of $\mathcal{F}_1, \mathcal{F}_2, \dots, \mathcal{F}_t$, respectively. Then
 $|\mathcal{F}_1|\geq {n-1\choose k_1-1}$ and $|\mathcal{F}_2|\geq {n-1\choose k_2-1}$, equivalently, $\{1, n-k_1+2, \dots, n\} \prec I_1$ and $\{1, n-k_2+2, \dots, n\} \prec I_2$.
\end{lemma}

\begin{proposition}\label{7-18-1}
Let $k_1\geq k_2\geq\dots\geq k_t\geq 2$ and $k_1+k_3\leq n <k_1+k_2$.
Suppose that $(\mathcal{F}_1, \dots, \mathcal{F}_t)$ is an L-initial extremal $(n, k_1, \dots, k_t)$-cross intersecting system with  IDs $I_1, I_2, \dots, I_t$  of $\mathcal{F}_1, \mathcal{F}_2, \dots, \mathcal{F}_t$, respectively. Then the following properties holds:
\begin{itemize}
\item[(i)] $I_2$ is the corresponding $k_2$-set of $I_1$;
\item[(ii)] for each $i\in [3, t]$, $I_i$ is the $k_i$-partner of $I_1$.
%\item[(iii)] $(I_1, I_3)$ is maximal.
\end{itemize}
\end{proposition}

\begin{proof}
By Lemma \ref{>}, we have $\{1, n-k_1+2, \dots, n\} \prec I_1$ and $\{1, n-k_2+2, \dots, n\} \prec I_2$.
For each $i\in [3, t]$,  let $P_i, Q_i$ be the $k_i$-partners of $I_1, I_2$, respectively.

First, we consider the case: $I_2\prec I_1$.
Let $i\in [3, t]$.
Since $n\geq k_1+k_i\geq k_2+k_i$ and $I_2\prec I_1$,
 by  Fact \ref{fact2.7}, $P_i\prec Q_i$. Note that $\mathcal{F}_1$ and $\mathcal{F}_i$ are cross intersecting and  $n\geq k_1+k_i$. By Fact \ref{f2.8}, $P_i$ is maximal to $I_1$. Therefore, $I_i\prec P_i$.
 By Remark \ref{remark2.20+}, $(\mathcal{F}_1, \mathcal{F}_2, \mathcal{L}(P_3, k_3), \dots, \mathcal{L}(P_t, k_t))$ is a cross intersecting system.
 Since
 $(\mathcal{F}_1, \dots, \mathcal{F}_t)$ is extremal, $I_i= P_1$, i.e., $I_i$ is $k_i$-partner of $I_1$, we get (ii). Therefore, $I_2$ is lexicographically maximal without exceeding $I_1$. Thus, $I_2$ is the corresponding $k_2$-set of $I_1$, we get (i).

Next, we consider the case: $I_1\prec I_2$.
Let $i\in [3, t]$.
Using a similar argument as above, we obtain that $I_i$ is $k_i$-partner of $I_2$.
By Fact \ref{fact2.7}, to prove Proposition \ref{7-18-1}, it suffices to prove that $I_1$ is the $k_1$-parity of $I_2$. To see this, we need the following claim.
\begin{claim}\label{7-18-2}
$(I_2, I_3)$ is maximal.
\end{claim}

\begin{proof}
Suppose on the contrary that  $(I_2, I_3)$ is not maximal.
Let $I'_2$ be the $k_2$-partner of $I_3$.
Note that $I_3$ is the $k_3$-partner of $I_2$.
Since $(I_2, I_3)$ is not maximal, in view of Fact \ref{f2.10},
$I_2\precneqq I'_2$ and $(I'_2, I_3)$ is maximal. So $I_2$ and $I'_2$ share the  same $k_3$-partner $I_3$. Consequently, they also share the  same $k_i$-partner $I_i$ for each $i\in [3, t]$.
Thus $(\mathcal{F}_1, \mathcal{L}(I'_2, k_2), \mathcal{F}_3, \dots, \mathcal{F}_t)$ is cross intersecting. However, $\mathcal{F}_2 \subsetneqq \mathcal{L}(I'_2, k_2)$, a contradiction to the fact that $(\mathcal{F}_1, \dots, \mathcal{F}_t)$ is extremal. This completes the proof of Claim \ref{7-18-2}.
\end{proof}

Let us continue the proof of Proposition \ref{7-18-1}.
Since $n\geq k_1+k_3$ and $k_1\geq k_2$, combining Proposition \ref{prop2.9}
and Claim \ref{7-18-2}, $I_2$ has $k_1$-parity. Since $I_1\prec I_2$ and $(\mathcal{F}_1, \dots, \mathcal{F}_t)$ is extremal, $I_1$ is the $k_1$-parity of $I_2$. This completes the proof of Proposition \ref{7-18-1}.
\end{proof}

\end{document}
